% !TeX root = ../../Book1.tex
% 纲要标题：### 6.3 由剖分到边识别多边形
% 纲要位置：几何学计划纲要.md，第 336 行。
% 正文状态：初稿；尚未正式审核。

% 纲要内容开始
% 
% * 生成树切割和单个多边形
% * 边配对及边界记录
% * 切割粘合操作的合法性
% 
% 纲要内容结束

\section{由剖分到边识别多边形}
\label{g1:sec:0603}

三角剖分把曲面化成有限资料，但三角形的个数通常与曲面的复杂程度相差很远。一个球面可以有成千上万个三角形，却不应需要成千上万个分类参数。本节先沿一棵对偶生成树把三角形展开成单个多边形，再把原来的顶点树吸收到一张圆盘中。留下的有限带条才记录曲面上不能靠改变分割消去的粘合。

\subsection{生成树切割和单个多边形}
\label{g1:subsec:0603:01}

固定紧连通曲面 \(M\) 的有限三角剖分，其存在性由\cref{g1:thm:surface-triangulation}保证。必要时作一次细分，以便在每个顶点周围看到由三角形组成的扇：内点的扇环绕一整圈，边界点的扇只绕半圈。原剖分的顶点与边组成一个有限图，称为其一维骨架。另放一个顶点在每个三角形内部；对每条非边界边，把它两侧的三角形中心经边的中点连接，得到\term[duioutu]{对偶图}[dual graph]。这些连线实际画在三角形中，只有在指定的边中点才穿过原骨架。

\begin{lemma}[对偶连通性]\label{g1:lem:surface-dual-connected}
	紧连通曲面的任一有限三角剖分，其对偶图连通。若曲面无边界，且 \(T\) 为一维骨架的一棵生成树，则删除与 \(T\) 的边相交的对偶边以后，对偶图仍连通。
\end{lemma}
\begin{proof}
	把经非边界边相邻的三角形归入同一类。同一内边两侧当然同类。在一个顶点周围，三角形按圆周或区间顺次排列，相邻者共享一条经过该顶点的非边界边，所以全部同类。于是各类三角形的并不但闭，而且在顶点、边内点和面内点处都有包含于该并的小邻域，因而也开。曲面连通，只有一类，证明第一项。

	设曲面无边界，删除指定对偶边后却出现两个以上分支。取其中一个分支的全部三角形为集合 \(\mathcal A\)。将一侧在 \(\mathcal A\)、另一侧不在 \(\mathcal A\) 的原边收集成子图 \(L\)。原对偶图连通，故 \(L\) 非空；而未被删除的对偶边不能连接两个新分支，故 \(L\subseteq T\)。

	在每个原顶点周围沿一圈看相邻三角形，进入 \(\mathcal A\) 与离开 \(\mathcal A\) 的次数相同。因此 \(L\) 在每个顶点的度数都是偶数。另一方面，树的任意非空边子图是森林，其中一个含边分支的最长简单边路有度数为一的端点。这与所有度数为偶数矛盾，故删边后的对偶图仍连通。
\end{proof}

\begin{theorem}[单个多边形表示]\label{g1:thm:surface-polygon}
	每个紧连通拓扑曲面都同胚于一张闭多边形按部分边的成对识别所得的商。闭曲面的每条边都参与且只参与一次配对；带边界时，未配对边的像组成曲面的边界。每一对边的识别均可在选定边参数后取为线性识别。
\end{theorem}
\begin{proof}
	选择对偶图的一棵生成树 \(T^*\)。先把各三角形当作彼此独立的副本，只沿 \(T^*\) 所穿过的原边粘合，不作其余任何识别。根三角形是一张圆盘。按树离根的距离逐个加入其他三角形，每次新三角形只沿一条完整边接到已有圆盘的一条边上；顶点和其他边仍保留各自副本。这样的粘合仍是圆盘：两张三角形或多边形沿边拼成的区域可把共同边直化，再在两侧使用半圆盘坐标；亦可把抽象的边界顶点依序放在一个凸多边形上，对角线连接后逐三角形定义分片线性同胚。归纳得到一张闭多边形 \(P\)。

	未穿过 \(T^*\) 的每条内边现在在 \(\partial P\) 上出现两次，按原边参数重新配对；原边界边只出现一次，不配对。顶点识别由这些边识别自动产生，因为同一原顶点周围的三角形扇是连通的。于是原三角形到 \(M\) 的映射下降为 \(P/{\sim}\to M\) 的连续双射。源空间紧，目标 Hausdorff，故由\cref{g1:thm:compact-bijection}为同胚。面内点不参与剩余识别，配对边内点有两张半圆盘，未配对边内点只有一张半圆盘；在顶点处，原有的圆形或区间形扇被恢复。因此未配对边的像恰是 \(\partial M\)。
\end{proof}

这一步只用了对偶树。一张单面多边形的各个顶点仍可能代表曲面中的不同点。下一步同时使用原骨架的一棵树，目的是把顶点附近的树状部分并成一张圆盘，而不是把一条任意闭路压成一点。

\begin{proposition}[一个顶点圆盘与有限带条]\label{g1:prop:one-vertex-band-model}
	设 \(M\) 为紧连通闭拓扑曲面。存在一张闭圆盘 \(D\)、有限条矩形带条和一张面圆盘 \(F\)，使 \(M\) 可按如下方式构成：每条带条的两个短边分别粘到 \(\partial D\) 上互不相交的小区间，长边保持为边界；\(D\) 与全部带条组成的曲面 \(R\) 恰有一个边界圆周，再沿此圆周粘上 \(F\)。若没有带条，则 \(M\cong S^2\)。
\end{proposition}
\begin{proof}
	在一维骨架中选生成树 \(T\)。由上引理，可在不穿过 \(T\) 的对偶图中选生成树 \(T^*\)，故两棵嵌入的树不相交。先围绕各原顶点取互不相交的小多边形圆盘，沿每条原边取窄矩形，在每个三角形中央留下一个小圆盘。这些选择可在有限三角形的重心细分中同时完成：顶点小片、边的窄条和面中央小片只沿各自边界相接，覆盖整张曲面。

	把属于 \(T\) 的边条加入各顶点小圆盘。每加一条树边，它连接的是此前不同的两个树分支，所以只是用一个矩形把两张圆盘接起来；所得仍是圆盘。按树归纳，所有顶点片及树边条合成一张圆盘 \(D\)。同样，把 \(T^*\) 所穿边处的横向窄条与各面中央圆盘合并，得到另一张圆盘 \(F\)。这两种合并不冲突，因为 \(T^*\) 没有穿过 \(T\)。

	尚余的区域恰对应既不在 \(T\) 中、也不被 \(T^*\) 穿过的原边。沿每条这样的边保留矩形带条，两个短边都接到 \(D\) 上，两个长边都接到面圆盘 \(F\) 的边界上。有限个多边形角可以在各原三角形内重新参数化成直角，不改变粘合。故 \(D\) 加全部带条的补部就是 \(\operatorname{Int}F\)，其边界为 \(\partial F\) 这一条圆周，得到所需描述。

	无带条时是两张圆盘沿整个边界粘合。任一边界粘合同胚 \(h:S^1\to S^1\) 可径向延拓为 \(ru\mapsto r h(u)\) 的圆盘同胚，故可把粘合改成恒等粘合；两盘分别放在上下半球便得到 \(S^2\)。
\end{proof}

在这个构造中，树并未被直接商成一个点。我们把树的一个盘状邻域当作新的顶点圆盘，所有改变都在已有曲面内进行。这也解释了为何带条的扭转仍须记录：树可以吸收，带条的两端以何种方向接到圆盘却影响整张曲面的定向与拓扑。

\subsection{边配对及边界记录}
\label{g1:subsec:0603:02}

给每对待识别边规定同一名称和箭头，沿多边形边界选定一个行走方向。顺着箭头读 \(a\)，逆着箭头读 \(a^{-1}\)，便得到循环的\term[bianzi]{边字}[boundary word]。闭曲面中每个字母恰出现两次，符号可能相同，也可能相反。未配对边另作标记，不与任何其他边识别。边字记录的是边的参数配对；顶点的等价类也必须由箭头逐项追踪，不能任意规定所有角为同一点。

\begin{example}[两种四边形]\label{g1:ex:polygon-vertex-counts}
	边字 \(aba^{-1}b^{-1}\) 的四个角属于同一等价类，给出环面。边字 \(aa^{-1}bb^{-1}\) 却有三个顶点等价类：依次记四角为 \(v_0,v_1,v_2,v_3\)，\(a\) 配对只令 \(v_0=v_2\)，\(b\) 配对也只令 \(v_2=v_0\)，而 \(v_1,v_3\) 各自保留。这时一维骨架是一棵有两条边的树，顶点盘加带模型没有剩余带条，所以得到球面。
\end{example}

两字在相应群关系中的表现很不同，但这里先按边识别构造空间，不能先对字作群的自由约化，再宣称所得商空间不变。允许的空间操作将在下个小节说明。

\begin{proposition}[边界领圈]\label{g1:prop:boundary-collar}
	设 \(M\) 为紧拓扑曲面。它的边界有一个\term[lingquan]{领圈}[collar]：存在嵌入
	\[
		c:\partial M\times[0,1]\longrightarrow M,\qquad c(x,0)=x,
	\]
	使 \(c(\partial M\times[0,1))\) 是 \(\partial M\) 的开邻域。不同边界分支的领圈可以取为互不相交。
\end{proposition}
\begin{proof}
	取含边界子复形的有限三角剖分。沿一条边界边的内部，在邻接三角形中取一条窄矩形；在边界顶点处，邻接三角形形成半圆盘状的有限扇。在这个扇中以一条靠近顶点的折线连接相邻两条边上的截点，折线与边界上的两小段围成一个圆盘。把它参数化成矩形，使下边对应原边界小段，左右两边接到上述边条，顶边对应内侧折线。这样的参数化存在：给四段边界指定对应，利用平面圆盘的边界延拓；也可在此有限扇中逐个三角形线性构造。

	于是每条边界圆周旁得到有限个矩形首尾相接，每个矩形的下边都在原边界上。相邻矩形在横截区间上的粘合固定下端和上端，所以其区间参数是递增同胚。依次改变各矩形的参数，把除最后一次之外的粘合都改为恒等；最后一条缝上的粘合是某个递增同胚 \(h:[0,1]\to[0,1]\)。它固定两端，且
	\[
		h_s(t)=(1-s)t+s h(t)\qquad(0\le s\le1)
	\]
	对每个 \(s\) 仍为递增同胚。从这个首尾粘合的矩形到标准环带的映射
	\[
		[s,t]\longmapsto(e^{2\pi i s},h_s(t))
	\]
	在 \((1,t)\sim(0,h(t))\) 处相容，是连续双射，因紧到 Hausdorff 而为同胚。这表明沿边界没有横向反转的可能：反转会把属于原边界的下端送到内侧上端。

	所得闭环带的外边是原边界，去掉内侧边后是相对开邻域；各块中的矩形坐标也直接核验这一点。边界只有有限个分支，且相互为不交紧集，先缩小所选扇片和边条即可使全部领圈不交。
\end{proof}

\begin{proposition}[封住边界]\label{g1:prop:cap-boundary}
	设 \(M\) 为有 \(b\) 个边界圆周的紧连通拓扑曲面。沿每个边界圆周用任意同胚粘入一张闭圆盘，所得空间 \(\widehat M\) 是紧连通闭曲面，且 \(M\) 同胚于 \(\widehat M\) 去掉这 \(b\) 张圆盘的内部。所得同胚类型不依赖封口同胚的选择。
\end{proposition}
\begin{proof}
	用上命题的领圈，在每条粘合缝附近使用圆周参数与法向参数：原曲面一侧的法向参数非负，所添圆盘薄环带一侧的参数非正。拼合以后成为圆周乘一个开区间，所以缝上各点都有平面圆盘邻域。新圆盘中心和原曲面内点都有原来的圆盘图，因而没有边界点。

	空间由有限个紧空间粘合而紧，并且连通。它是 Hausdorff：不同缝上的点可取不交薄领圈；同一缝的不同点取不交圆弧乘薄区间；不在缝上的点可在各自原片内取避开缝的邻域，或先把薄领圈缩小后分开。这些邻域也是商拓扑中的开集。有限组来自原曲面和添入圆盘的可数坐标基给出第二可数性，因此所得确为闭拓扑曲面。

	若改变某个封口同胚，两种选择之差为圆周同胚。把它径向延拓到所添圆盘，原曲面上保持恒等，即得到两个商空间之间的同胚。这也保持原曲面的副本，故不仅抽象空间同胚，原曲面及封口位置也相容。
\end{proof}

所以有边界情形可以先封盘化为闭曲面，再追踪取下了多少张盘。但“同样数目的盘可以移到同样的位置”仍需证明，下一节恢复边界时会补上这个位置无关性，而不把它藏在图形直觉中。

\subsection{切割粘合操作的合法性}
\label{g1:subsec:0603:03}

边字只是表示方法，下面的变化应当被理解为同一块材料重新切开和接合。若没有把边界参数与剩余识别一并带走，即使变换后的群展示同构，也不能据此断言两张曲面同胚。

\begin{lemma}[圆盘的延拓与相对边界同痕]\label{g1:lem:disk-extension}
	任意圆周同胚 \(h:S^1\to S^1\) 都延拓为圆盘同胚。若圆盘同胚 \(f:D^2\to D^2\) 在边界上恒等，则存在从恒等映射到 \(f\) 的连续同胚族，且族中每个同胚都在边界上恒等。
\end{lemma}
\begin{proof}
	第一项用 \(H(ru)=r h(u)\)，其中 \(u\in S^1,0<r\le1\)，并令 \(H(0)=0\)。它与以 \(h^{-1}\) 定义的同类公式互逆；在原点处由像点的模等于 \(r\) 得到连续性。

	第二项对 \(0<t\le1\) 定义
	\[
		F_t(x)=
		\begin{cases}
			t f(x/t),&\lVert x\rVert\le t,\\
			x,&t\le\lVert x\rVert\le1,
		\end{cases}
		\qquad F_0(x)=x.
	\]
	在 \(\lVert x\rVert=t\) 处两式一致，因为 \(f\) 固定单位圆。每个 \(F_t\) 的逆由相同公式中的 \(f^{-1}\) 给出，故为同胚。\(t>0\) 时分片粘合保证联合连续；当 \(t\to0\) 时，所有移动只发生在半径 \(t\) 的圆盘内，且 \(\lVert F_t(x)-x\rVert\le2t\)，故在 \(t=0\) 也联合连续。最后 \(F_1=f\)，并始终固定单位圆。
\end{proof}

上述连续同胚族也称为\term[tonghen]{同痕}[isotopy]；把它在圆盘外接上恒等映射，就得到\cref{g1:def:ambient-isotopy}的环境同痕。它比一般同伦严格：每一时刻都必须可逆。本章用它调整切割位置时，曲面同胚类型始终保留。

\begin{proposition}[多边形表示的基本变化]\label{g1:prop:polygon-moves}
	对已表示拓扑曲面的边配对多边形，下列操作保持所得曲面的同胚类型：
	\begin{enumerate}
		\item 同时重命名一对边、反向选择该边的箭头，或改变边界字的起读位置；也可反向走完整个多边形边界，此时次序和所有指数同时反转。
		\item 在一对边上加入相应的分点，或删去这种成对分点。
		\item 沿多边形内一条连接不同边界点的适当嵌入简单弧切成两张圆盘，给两份切缝加入相应识别；随后可选择其他一对完整边将两盘重新合为一盘，所有其余配对按原参数保留。
	\end{enumerate}
	最后一种操作中的两张盘允许先各自换用圆盘坐标或翻面，但必须同时带走边界箭头及其参数。
\end{proposition}
\begin{proof}
	前两项只改变同一商映射的参数名称或分段，等价关系未变。第三项由\cref{g1:lem:proper-arc-straightening}将切弧化为圆盘直径，因此切开后确是两盘，按原参数重新接上切缝便逐点找回原盘。若改先粘另一对边，抽象的两盘沿一段完整边拼合仍是圆盘；其边界各段继承原识别。无论先作哪次粘合，最终对两盘不交并所加的等价关系完全相同，所以两种分步商都连续双射到同一个原曲面。紧到 Hausdorff 的判据将这些双射提升为同胚。

	圆盘坐标的改变由上引理将边界参数延拓到内部；翻面也是圆盘同胚，只是改变边界行走方向。因此只要同步改变箭头及边界次序，最后的等价关系仍相同。若把箭头保留而单独翻面，则可能改变识别，这不属于所述合法操作。
\end{proof}

一个常用的特例是消去边字中的相邻 \(aa^{-1}\)。这两边在商图中形成一条从其余图伸出的末端边；围住这条边及其末端顶点的小邻域是圆盘，把它并入相邻面就去掉了这个“树枝”。若还有其他边，得到余下的边字；若整字只有 \(aa^{-1}\)，则商空间由两个盘组成，是球面，不能把空字误读为空空间。同号的 \(aa\) 不属于这种消去，它记录了不同的粘合。

\ProseParagraphBreak

下一节将反复使用这些圆盘与带条操作抽出手柄和交叉帽。证明的终止性由带条数或边对数的严格减少保证；Euler 数稍后用于区分最终模型，不能提前拿来替代存在性中的切割论证。
